Here and throughout this paper primes denote radial derivatives, and commas followed
by $t$'s denote time derivatives: $X' \equiv \partial X/\partial r$, $X'' \equiv \partial^{2}X/\partial r^{2}$,
$X_{,t} \equiv \partial X/\partial t$, and $X_{,tt} \equiv \partial^{2}X/\partial t^{2}$.

In deriving equations (8) and (9) we have assumed at several points that $l \neq 0$ and
$l \neq 1$; and \textit{we shall maintain this assumption throughout the remainder of the paper}. This
assumption is not a serious limitation on our analysis, since we are interested primarily
in the emission of gravitational waves by pulsating stars, and gravitational waves must
have $l \geq 2$.

The structure of the equations of motion (8) and (9) reflects the fact that even-parity
motions have three degrees of freedom---two, $W$ and $V$, associated with the motion of
the fluid; and one, $K$, associated with the gravitational waves.\footnote{If one specifies at an
initial moment of coordinate time the radial distributions of $W$, $V$, $K$ and $W_{,t}$, $V_{,t}$, $K_{,t}$,
then one can use the three propagation equations (9)---plus the initial value equation
(8a) with the boundary condition
\begin{equation}
  H_0 = K \quad \text{at } r = 0
\label{eq:10}
\end{equation}
(cf.\ eq.\ [15a])---to propagate $W$, $V$, $K$ and $K_{,t}$. At any moment of coordinate
time the functions $H_0$, $H_1$, and $H_2$, which are needed to fix completely the perturbed
geometry of spacetime, are determined from $W$, $V$, $K$, $V_{,t}$, $K_{,t}$ by the initial-value
equations (8).

The dynamical evolution problem (eqs.\ [9]) would be much simpler if one could solve
the initial-value equations (8a) for $H_2$ as a linear combination of $K'$, $K$, $K_{,t}$, $W'$, $W$,
$V''$, $V'$, $V$, and then use that solution, plus equation (8a) itself, to eliminate $H_0$ and $H_2$
from the dynamical equations (9). Unfortunately, equations (8a) cannot be so inverted
and, therefore, equation (8a) is needed along with equations (9) to fix the dynamical
propagation of the three degrees of freedom. The only way one can incorporate equation
(8a) into the propagation equations (9) directly is by increasing equations (9) from
second-order partial differential equations to third-order partial differential equations. We
prefer not to do this.}

\section{EVEN-PARITY NORMAL MODES OF PULSATION}

The study of the even-parity pulsations of the equilibrium configuration is simplified
considerably by analyzing the pulsations into normal modes. For pulsations the only pulsations
one expects to occur in nature are pulsations which generate outgoing gravitational
waves; and these are damped primarily (but not entirely) in models in which the
gravitational radiation is purely outgoing. Because of radiation damping, such normal
modes will not have real frequencies and eigenvalues; rather, their eigenfrequencies
and frequencies will be complex. However, assuming completeness of the set of
complex eigenfunctions with outgoing waves, any real pulsation with purely outgoing
radiation can be expressed as a linear combination of the complex normal modes.

In this section we shall discuss the eigenvalue equations and the boundary conditions
which determine the normal modes of even parity---including real normal modes (standing
gravitational waves), as well as complex normal modes with various mixtures of
outgoing and incoming gravitational radiation. All discussion of the relationship between the normal modes and the real pulsations of an equilibrium configuration will
be delayed until \S~IV.
